%!TEX root = ../atomic_lipics.tex
% LTeX: language=en
%
\section{Background}
%
\subsection{Ideals and polynomial rings}
%
Let $\K$ be an arbitrary field and $X$ be a set of variables.
We use $\poly{\K}{X}$ denote the polynomial ring over $\K$ in these variables.

A subset $\idl[I] \subseteq \poly{\K}{X}$ is an \intro{ideal} if
\begin{enumerate}
\item $0 \in \idl[I]$,
\item for every $f,g$ in $\idl[I]$, their sum $f + g$ is also in $\idl[I]$, and
\item for every $f\in\idl[I]$ and $h\in\poly{\K}{X}$, their product $f \cdot h$ is also in $\idl[I]$.
\end{enumerate}
%
\begin{example}\label{eg:I0}
Let $\idl[I]_0^X$ be the set of polynomials whose constant term is $0$.
For illustration,
the polynomial $x^2 + y$ is in $\idl[I]_0^{\{x,y\}}$ and the polynomial $xy + 2$ is not.
We leave it to the reader to check that $\idl[I]_0^X$ is indeed an ideal. 
\end{example}
%
\arka{possible repetition}.
%
\begin{example}\label{eg:I1}
Let $\idl[I]_1^X$ be the set of polynomials whose coefficients add up to $1$.
For illustration,
the polynomial $x^2 - 2y + 1$ is in $\idl[I]_0^{\{x,y\}}$ and the polynomial $xy + y^2$ is not. 
\end{example}
%
For a subset $F \subseteq \poly{\K}{X}$,
the ideal \intro{$\IdlGen{F}$} generated by $H$ is the smallest ideal containing $F$.
%
\begin{remark}
The ideal $\IdlGen{F}$ generated by $F\subseteq \poly{\K}{X}$ can equivalently defined as the set of all polynomials that can be written as
$
\sum_{i = 1}^n h_i \cdot f_i
$
for $h_i\in\poly{\K}{X}$ and $f_i \in F$.
\end{remark}
%
\begin{example}
The ideals $\idl[I]_0$ and $\idl[I]_1$ defined in \Cref{eg:I0,eg:I1} are generated respectively by the sets $\setof{x}{x\in X}$ and $\setof{x - 1}{x\in X}$.
\end{example}
%
\arka{Ideal variety correspondence. To motivate the need of computing sum,
intersection and product?
Depending on space}
%
\subsection{Polynomial rings with finitely many variables}
%
An ideal $\idl[I]$ is called \intro{finitely generated} if it is generated by a finite set of polynomials.
For instance, when $X$ is finite, the ideals $\idl[I]_0^X$ and $\idl[I]_1^X$ are finitely generated.
In fact,
%
\begin{theorem}[Hilbert's basis theorem]
If $X$ is finite,
then every ideal in $\poly{\K}{X}$ is finitely generated.
\end{theorem}
%
In the rest of this subsection we will assume that the set of variables $X$ is finite.

Hilbert's basis theorem implies that ideals in $\poly{\K}{X}$ can be finitely represented (assuming elements of the field $\K$ can also be finitely represented)and given as input to algorithms .
One of the most important decision problem regarding ideals is the \intro{ideal membership problem}:
decide whether a given polynomial $f\in\poly{\K}{X}$ is in the ideal generated by a given finite set of polynomials $F\subseteq\poly{\K}{X}$.
 



\arka{repeated division}

\arka{sufficient but not necessary}

\arka{show with example how to compute THE generator for univariate}

\arka{multivariate}

\arka{term order}

\arka{reverse lex order as an example}

\arka{division}

\arka{grobner basis}

\arka{why a set may not be grobner basis and motivate the Buchberger's condition}

\arka{buchberger}

%
\begin{enumerate}
\item Ideals
\item Divisibility and reduction
\item Ideal membership
\item Gr\"{o}bner basis 
\item Buchberger's algorithm
\end{enumerate}